\section{Checkpoint 4: train the circuit angles}\label{sec:training}
\subsection{What the optimizer minimizes}
For angles $\boldsymbol\beta=(\beta_1,\ldots,\beta_p)$ and
$\boldsymbol\gamma=(\gamma_1,\ldots,\gamma_p)$, let
\begin{equation}
p_{\boldsymbol\beta,\boldsymbol\gamma}(\boldsymbol{x})
=\left|\langle\boldsymbol{x}\mid
\psi_p(\boldsymbol\beta,\boldsymbol\gamma)\rangle\right|^2.
\end{equation}
The bra $\bra{\boldsymbol{x}}$ is the conjugate transpose of the corresponding
ket. The inner product extracts that outcome's amplitude. Squaring its magnitude
gives the probability. Appendix~\ref{app:math} explains the notation in detail.

The expected QUBO score is the probability-weighted average
\begin{equation}\label{eq:expected-cost}
F(\boldsymbol\beta,\boldsymbol\gamma)
=\sum_{\boldsymbol{x}}p_{\boldsymbol\beta,\boldsymbol\gamma}(\boldsymbol{x})
f(\boldsymbol{x})
=s\bra{\psi_p}\widetilde H_C\ket{\psi_p}+C_I.
\end{equation}
The sum includes invalid booking strings and their penalties. A circuit can
occasionally produce an excellent valid schedule while still having a poor
average because many other measurements would be invalid.

The implementation passes the normalized expectation to COBYLA:
\begin{lstlisting}[language=Python]
pub = (circuit, normalized_cost, angles)
energy = float(
    estimator.run([pub], precision=0).result()[0].data.evs
)
expected_score = energy * scale + offset
\end{lstlisting}
\code{precision=0} asks the local Aer estimator for the ideal expectation,
computed to numerical precision without finite-shot sampling noise\cite{aer}.
The quantum state is still simulated on a classical computer. Final measurements
are sampled separately; \code{--shots} does not change the training precision.

\subsection{What COBYLA changes}
COBYLA is a derivative-free classical numerical optimizer\cite{cobyla}. It
proposes circuit angles, receives an expectation, and proposes new angles.
It does not choose booking bits directly and does not call an exhaustive
scheduling solver in the quantum path.
\begin{equation*}
\text{angles}\ \longrightarrow\ \text{Aer circuit evaluation}
\ \longrightarrow\ \text{expectation}\ \longrightarrow\ \text{new angles}.
\end{equation*}
The implementation uses \code{maxiter} as an objective-evaluation budget per
restart, initial step size \code{rhobeg=0.5}, and trust-region tolerance
\code{tol=1e-4}. The CLI requires at least $2p+2$ evaluations. Convergence may
occur before the budget is exhausted; no global optimum is certified.

Different restarts use different initial angles. The seed initializes a local
random generator, and the best parameter vector actually evaluated across all
restarts is retained. With $p=2$, Qiskit's parameter order for this circuit is
\begin{equation}
(\beta_1,\beta_2,\gamma_1,\gamma_2).
\end{equation}
In Python these are \code{beta[0]}, \code{beta[1]}, \code{gamma[0]}, and
\code{gamma[1]}. This is the binding order, not the execution order of the gates.
Use the returned \code{parameter_order} instead of guessing if the circuit changes.

\subsection{A concrete training trace}
For the current code and default seed, the three restarts used 92, 146, and 78
evaluations respectively, for 316 total. The best expected score observed was
approximately
\begin{equation}
F_{\mathrm{best}}\approx73.88034.
\end{equation}
This is not the final selected schedule's score of $46.30$. Equation~\eqref{eq:expected-cost}
explains the difference: the expectation averages the whole distribution.
The pure uniform starting state would have mean score $C_I=227.7125$, because
all nonidentity Z and ZZ expectations vanish. A restart's first evaluated
angles are random, nonzero angles; its initial expectation need not equal that
uniform-state mean.

\section{Checkpoint 5: measure, decode, and check}
\subsection{Sampling the trained circuit}
After binding the best angles, append measurements and execute the circuit:
\begin{lstlisting}[language=Python]
measured = circuit.assign_parameters(best_angles)
measured.measure_all()
counts = simulator.run(
    measured, shots=4096, seed_simulator=7
).result().get_counts()
\end{lstlisting}
Each dictionary entry is a displayed bit string and the number of shots that
produced it. The entries count circuit outcomes, not optimizer evaluations.

\subsection{Qiskit bit order}
Logical booking vectors in this guide use $(x_0,x_1,x_2,x_3,x_4)$. Qiskit displays
the classical bits in reverse index order:
\begin{equation}
\text{displayed string}=x_4x_3x_2x_1x_0.
\end{equation}
Thus the displayed string \code{10010} corresponds to
\begin{equation}
(x_0,x_1,x_2,x_3,x_4)=(0,1,0,0,1).
\end{equation}
The decoder explicitly reverses the string:
\begin{lstlisting}[language=Python]
bits = [int(bit) for bit in bitstring[::-1]]
\end{lstlisting}
This matters for every mapping from measurement to a real booking. It is not
just a cosmetic choice in how a histogram is printed.

\subsection{Validity and sampling metrics}
The function \code{decode_bits()} checks the selected bookings independently of
the penalty coefficients: exactly one booking per car and no shared station/slot.
Let $N$ be the final shot count, $n_{\boldsymbol{x}}$ the count for a string, and
$\mathcal F$ the set of feasible strings. The empirical valid fraction is
\begin{equation}\label{eq:valid-rate}
\widehat p_{\mathrm{valid}}=\frac1N
\sum_{\boldsymbol{x}\in\mathcal F}n_{\boldsymbol{x}}.
\end{equation}
This counts shots, not just distinct valid strings. A valid string observed
1,000 times contributes 1,000 to the numerator.

\begin{center}
\begin{tabular}{lrrl}
\toprule
Displayed string & Shots & QUBO score & Interpretation\\
\midrule
\code{00001} & 719 & 79.70 & A only; B unassigned\\
\code{10010} & 663 & 46.30 & A North late; B North early\\
\code{10100} & 629 & 52.65 & A South early; B North early\\
\code{11000} & 629 & 56.65 & A South late; B North early\\
\code{10001} & 250 & 99.35 & Both cars at North early\\
\bottomrule
\end{tabular}
\end{center}
These are selected entries from the default run, not the complete histogram.
The most frequent outcome was invalid. The three valid strings contributed
$663+629+629=1921$ shots, giving
\begin{equation}
\widehat p_{\mathrm{valid}}=\frac{1921}{4096}\approx0.468994=46.90\%.
\end{equation}
For reference, the uniform distribution has only three feasible strings among
32 and hence valid probability $3/32=9.375\%$. This comparison evaluates the
distribution; it is not evidence of computational speedup over classical solvers.

\section{Checkpoint 6: choose a sampled schedule and report it}
The program selects the lowest booking score among feasible strings that were
actually measured. It does not manufacture an unobserved string, repair an
invalid sample, or invoke the reference exact solver as a fallback.

For the default run the selected logical vector was $(0,1,0,0,1)$:
\begin{align*}
\text{A}&:\ \text{North, reserved 09:00--10:00},\\
\text{B}&:\ \text{North, reserved 08:00--09:00},\\
\text{electricity bill}&=10.40+18.40=\text{Rs }28.80,\\
\text{waiting}&=65+5=70\text{ minutes},\\
\text{schedule score}&=28.80+0.25(70)=46.30.
\end{align*}
Actual charging finishes approximately at 09:56:44 for A and 08:50:11 for B.
Both ports remain reserved until the end of their assigned slots.

Successful results use \code{status="feasible"}, not \code{"optimal"}.
If no feasible string is measured, the result is \code{no_feasible_sample};
this does not prove that no feasible string exists. If preprocessing finds a
vehicle with no individual option, the CLI can report that definite obstruction
before building a circuit. These are different kinds of failure.

The 33 tests in the current snapshot cover physical feasibility, QUBO penalties,
all 32 basis-state energies, actual cost-gate phases, bit ordering, sampled-only
selection, small Aer runs, and global conflict cases. The tests may enumerate
tiny state spaces to verify correctness; the production quantum solver does not
use those enumerations to select a schedule.

\section{Experiments: what to improve and how to measure it}\label{sec:experiments}
Change one setting at a time and compare several seeds. Record the best feasible
sample score, valid fraction, best expected score, evaluation count, and runtime.
If a run fails to sample any valid schedule, keep that failure in the results.

\begin{center}
\begin{tabularx}{\textwidth}{@{}lX X@{}}
\toprule
Setting & Changes & Interpretation\\
\midrule
\code{--layers} & Number $p$ of cost/mixer pairs & More gates and $2p$ angles; more expressiveness does not guarantee better training.\\
\code{--iterations} & Evaluation budget per restart & More opportunities to optimize; convergence may occur earlier.\\
\code{--restarts} & Number of angle initializations & Helps explore different local solutions; increases total work.\\
\code{--shots} & Final measurement repetitions & Changes sampling coverage, not ideal training precision or angles.\\
\code{--seed} & Initialization and final sampling randomness & Tests sensitivity to favorable random choices.\\
\code{--penalty} & Constraint coefficient $P$ & Alters feasibility emphasis relative to operating costs.\\
\bottomrule
\end{tabularx}
\end{center}

\begin{lstlisting}
# Inspect the simplest cost/mixer circuit.
.venv/bin/python main.py --layers 1 --show-circuit

# Increase expressiveness and training effort together.
.venv/bin/python main.py --layers 3 --iterations 300 --restarts 5

# Increase final sampling coverage only.
.venv/bin/python main.py --shots 16384

# Repeat with a different initialization and sampling seed.
.venv/bin/python main.py --seed 21
\end{lstlisting}

The returned \code{history}, \code{angles}, and \code{parameter_order} make
learning-curve plots and angle comparisons possible. A later extension could
change the mixer, initialize from a useful state, simulate device noise, or use
finite-precision expectation estimates. These are experiments to implement;
they are not features already enabled by the current CLI.

Larger shot counts improve the probability of observing an existing rare useful
outcome. If its probability is $p_\star$ and shots are independent, the chance
of seeing it at least once in $N$ shots is
\begin{equation}
1-(1-p_\star)^N.
\end{equation}
Shots cannot reveal an outcome with exactly zero probability. The validity
fraction and training expectation tell us more than the best observed score alone.

\section{Boundaries, hardware access, and the browser explainer}
\subsection{What the current physical model assumes}
Each car needs one uninterrupted session within one nonoverlapping hour slot.
Each station has one port. Consumption and charging power are constant, charging
efficiency is 100\%, and the driver departs as soon as available. Actual charging
completion determines deadline compliance, but the reservation blocks the full
slot. Waiting is the time at the station before the slot starts; driving has no
separate objective penalty. Reserve is checked on arrival, not on a later onward
journey. These choices simplify the application and must be revisited when the
product requirements change.

\subsection{Simulation versus a physical quantum processor}
Qiskit builds the circuits; Aer simulates them locally. The IBM API key is not
used by this program. A physical IBM processor would require a separate Runtime
integration, backend selection, compilation to supported gates/connectivity,
and a training/measurement workflow compatible with that backend. An API key
alone does not change \code{AerSimulator()} into hardware execution.

An ideal statevector stores $2^n$ complex amplitudes. A compact Pauli operator
representation does not remove that simulator memory cost. The current
single-session allocation model is also approachable by strong classical
assignment methods. This project demonstrates a quantum optimization workflow;
it makes no claim that QAOA is necessary for this model or faster than classical
computing.

\subsection{The interactive parking illustration}
The separate browser illustration in \code{web/parking-anatomy/} shows the same
parking/charging example, waiting preferences, and a five-bit booking inspector.
Build it with
\begin{lstlisting}
node web/parking-anatomy/build.mjs
\end{lstlisting}
Then open \code{web/parking-anatomy/dist/index.html}. Its build exports sample
bookings from Python and can use a saved fixture. The browser computes 32
illustrative QUBO energies classically. It does not run QAOA or connect to IBM;
its exported baseline uses the reference classical solver. Run \code{main.py}
to execute the actual Qiskit circuit. The browser's own README contains its
preview instructions.

\subsection{Code-to-concept map}
\begin{center}
\begin{tabularx}{\textwidth}{@{}lX@{}}
\toprule
Function & Responsibility\\
\midrule
\code{sample_scenario} & Input dictionaries for cars, stations, slots, and travel.\\
\code{generate_options} & Physical feasibility, journey quantities, and booking scores.\\
\code{build_qubo} & Data-driven variables, exactly-one groups, and conflict penalties.\\
\code{qubo_value} & Evaluate a supplied bit vector; no optimization.\\
\code{qubo_to_ising} & Collect Z/ZZ coefficients, offset, and normalization.\\
\code{build_qaoa_circuit} & Explicit H, RZ, RZZ, and RX operations.\\
\code{solve_qaoa} & Train angles on Aer and measure the final circuit.\\
\code{decode_bits} & Independently check selected booking constraints.\\
\code{summarize_counts} & Reverse display order and select a feasible measured result.\\
\bottomrule
\end{tabularx}
\end{center}
The main quantum implementation is \code{src/quantum_scheduler.py}; ordinary
physical calculations are in \code{src/ev_scheduler.py}. The CLI is in
\code{main.py}.

\clearpage
